\section{Benchmark processes and validation}\label{sec:benchmark}

We dedicate this section to validating our LO implementation of the orbital and total angular-momentum projectors (eqs.~\eqref{eq:proj_oam} and \eqref{eq:proj_tam}), as well as the reshuffling routines {applied to eq.~\eqref{eq:partonic_xsec_explicit}.} For the projector implementation, in section~\ref{sec:Pwavetest} we compare our results for a single phase-space point (using the \textit{standalone mode}) and for the phase-space-integrated cross section against \helaconia~\cite{Shao:2012iz,Shao:2015vga}. Our tests involve processes up to three {\pwavetxt} quarkonia with both QCD and electroweak interactions included.
For the momentum reshuffling, we compare our algorithms against the non-reshuffling case in section~\ref{sec:momreshufftest}, obtained through an appropriate reparametrisation of the quark and quarkonium masses.
These benchmarks are essential to verify the correctness of our implementation and ensure the reliability of the applications discussed in sections~\ref{sec:quarkonium} and \ref{sec:leptionium}.

\subsection{Validation of the P-wave implementation}
\label{sec:Pwavetest}
With a slightly modified setup relative to the baseline configuration described in appendix~\ref{app:setup}, we benchmark the implementation of the new orbital and total angular momentum projectors. To match the parameter setup used in \helaconia{}, which treats the decay widths in internal propagators differently from \mgshort, we set the decay widths of the weak and Higgs bosons to zero, ${\Gamma_W = \Gamma_Z = \Gamma_H = 0}$. Moreover, we set all quarkonium masses to $M_{\cal Q}^{\phantom{2}}=m_{Q}^{\phantom{2}}+m_{\bar{Q}^\prime}$, where $m_{Q}^{\phantom{2}}$ and $m_{\bar{Q}^\prime}$ denote the masses of the heavy quark and antiquark constituents, respectively. This choice avoids momentum reshuffling, which is not implemented in \helaconia{}. Additionally, in the \textit{standalone mode}, we fix $\alpha_s=0.118$, while for integrated cross sections the value of $\alpha_s$ is obtained from {\tt LHAPDF6}~\cite{Buckley:2014ana} and evaluated at the renormalisation and factorisation scales defined in eq.~\eqref{eq:setup_mu}.

Our first comparison is performed at the level of the squared helicity amplitude, $|{\cal M}|^2$, averaged over the colours and helicities of the initial-state particles, summed over those of the final-state particles, and weighted by LDMEs. More precisely, these squared amplitudes correspond to the quantity obtained by multiplying the factor in front of the phase-space measure in eq.~\eqref{eq:partonic_xsec_explicit} by the corresponding LDMEs. For a $2 \to {\rm N}$ partonic process producing a set $\mathfrak{B}$ of bound states ${\cal B}[n]$, we obtain 
\begin{equation}
    \left|\mathcal{M}_{\mathrm{N},\mathfrak{B}}\right|^2 = \frac{1}{{\cal N}(\dot{r})}\!\left(\frac{1}{4k_1\cdot k_2} \frac{1}{\omega(\Ione)\omega(\Itwo)} \right)\! \prod_{{\cal B}[n] \in \mathfrak{B}} \!\left(\frac{1/(2\mu_{\cal B})}{(2J+1)N_{[C]}}\langle {\cal O}_n^{\rm B}\rangle\right)\!\sum_{\substack{\rm colour\\ \rm spin}} |  {\mathds A}^{(\mathrm{N},0)}(\dot{r})|^2\,.
\label{eq:me_square}
\end{equation}
Each squared amplitude is evaluated at a single phase-space point. While the {\it standalone mode} of \mgshort\ employs {\tt Rambo}~\cite{Kleiss1986RAMBO} by default, the generated phase-space points have limited numerical accuracy, leading to violations of the on-shell conditions in double-precision arithmetic. Therefore, we use an alternative implementation based on ref.~\cite{Byckling:1971vca}, which generates momenta with analytic precision. These momenta are then passed to both \helaconia\ and \mgshort\ in double precision. To avoid excessively boosted phase-space configurations, we set the partonic c.m.\@ energy for all processes to $\sqrt{\hat s} = \sum_i M_i + 20\,{\rm GeV}$, namely $20\,{\rm GeV}$ above threshold. 
The agreement between the two tools is quantified through the relative deviation, defined as
\begin{equation}
    \Delta_\mathrm{rel.} = 2\left|\dfrac{\left|\mathcal{M}_{N,\mathfrak{B}}\right|^2_{\mgshort} - \left|\mathcal{M}_{N,\mathfrak{B}}\right|^2_\helaconia}{\left|\mathcal{M}_{N,\mathfrak{B}}\right|^2_{\mgshort} + \left|\mathcal{M}_{N,\mathfrak{B}}\right|^2_\helaconia}\right|,
\label{eq:relative_ratio}
\end{equation}
where $|{\cal M}_{N,\mathfrak{B}}|^2_{\mgshort}$ and $|{\cal M}_{N,\mathfrak{B}}|^2_{\helaconia}$ denote the squared amplitudes as defined in eq.~\eqref{eq:me_square}, evaluated with \mgshort\ and \helaconia, respectively.
%%%%%%%%%%%%%%%%%%%%%%%%%%%%%%%%%%%%%%%%%%%%%%%%%%%%%%%%%%%%%%%%%%%%%%%%%%%
\begin{table}[H]
\centering
\renewcommand*{\arraystretch}{1.1}
\begin{tabular}[t]{lcc}
\toprule
\multirow{2}{*}{\textbf{process}}& \madgraph & \multirow{2}{*}{$\Delta_\mathrm{rel.}$}  \\
 & \helaconia  \\
\midrule
\multirow{2}{*}{$g g \to h_c\big[{}^1\pwave_{1}^{[1]}\big] + g u \bar{u}$} & $9.0256873979004725 \cdot 10^{-5}\,\mathrm{GeV}^{-4}$ & \multirow{2}{*}{$1.46 \cdot 10^{-14}$} \\
 & $9.0256873979006040 \cdot 10^{-5}\,\mathrm{GeV}^{-4}$ & \\ 
\hdashline

\multirow{2}{*}{$u \bar{u} \to \Upsilon\big[{}^3\pwave_{0}^{[8]}\big] + g g g$} & $1.1548805300001495 \cdot 10^{-5}\,\mathrm{GeV}^{-4}$ & \multirow{2}{*}{$1.38 \cdot 10^{-14}$} \\
 & $1.1548805300001654 \cdot 10^{-5}\,\mathrm{GeV}^{-4}$ & \\
\hdashline

{$g g \to \eta_c\big[{}^1\swave_{0}^{[8]}\big] + h_c\big[{}^1\pwave_{1}^{[1]}\big]$} & 
$4.7791765795384806 \cdot 10^{-16}\,\mathrm{GeV}^{-2}$ & \multirow{2}{*}{$9.06 \cdot 10^{-14}$} \\
\qquad\qquad $+ \Upsilon\big[{}^3\swave_{1}^{[1]}\big]$ & $4.7791765795389135 \cdot 10^{-16}\,\mathrm{GeV}^{-2}$ & \\
\hdashline

{$g \gamma \to \eta_b\big[{}^1\swave_{0}^{[1]}\big] + \eta_c\big[{}^1\swave_{0}^{[8]}\big]\ \textsuperscript{\ref{fnt:QED_flag}}$} & 
$1.0744719179339013 \cdot 10^{-13}\,\mathrm{GeV}^{-2}$ & \multirow{2}{*}{$2.94 \cdot 10^{-15}$} \\
\multirow{1}{*}{\qquad\qquad$ + \chi_c\big[{}^3\pwave_{2}^{[1]}\big]$} & 
$1.0744719179339045 \cdot 10^{-13}\,\mathrm{GeV}^{-2}$ & \\ 
\hdashline

{$u \bar{u} \to \chi_{c0}\!\big[{}^3\pwave_{0}^{[1]}\big] + \Upsilon\!\big[{}^3\pwave_{1}^{[8]}\big]$} & $2.9379230723426576 \cdot 10^{-16}\,\mathrm{GeV}^{-2}$ & \multirow{2}{*}{$5.71 \cdot 10^{-15}$} \\*
\qquad\qquad$+\eta_c\!\big[{}^1\pwave_{1}^{[8]}\big]$ & $2.9379230723426408 \cdot 10^{-16}\,\mathrm{GeV}^{-2}$ & \\
 \hdashline

{$g u \to J/\psi\big[{}^3\pwave_{1}^{[8]}\big] + \eta_b\big[{}^1\pwave_{1}^{[8]}\big]\ \textsuperscript{\ref{fnt:QED_flag}}$} & 
$1.0252237742951235 \cdot 10^{-18}\,\mathrm{GeV}^{-4}$ & \multirow{2}{*}{$1.10 \cdot 10^{-13}$} \\
\qquad\qquad$ + W^{+} d$ & $1.0252237742950112 \cdot 10^{-18}\,\mathrm{GeV}^{-4}$ & \\ 
\hdashline

\multirow{2}{*}{$g g \to \chi_b\big[{}^3\pwave_{1}^{[1]}\big] + Z H H \ \textsuperscript{\ref{fnt:QED_flag}}$} & $7.1410224887725591 \cdot 10^{-19}\,\mathrm{GeV}^{-4}$ & \multirow{2}{*}{$4.05 \cdot 10^{-16}$} \\
& $7.1410224887725562 \cdot 10^{-19}\,\mathrm{GeV}^{-4}$ & \\ 
\hdashline

\multirow{1}{*}{$ g g \to \chi_c\big[{}^3\pwave_{0}^{[1]}\big] + \chi_c\big[{}^3\pwave_{1}^{[1]}\big]$} & $9.8531189707367133 \cdot 10^{-13}\,\mathrm{GeV}^{-2}$ & \multirow{2}{*}{$6.76 \cdot 10^{-15}$} \\
\qquad\qquad$+ J/\psi\big[{}^3\swave_{1}^{[8]}\big]$ & $9.8531189707366466 \cdot 10^{-13}\,\mathrm{GeV}^{-2}$ & \\ 
\hdashline

\multirow{2}{*}{$g g \to J/\psi\big[{}^3\pwave_{2}^{[8]}\big] + \eta_c\big[{}^1\pwave_{1}^{[8]}\big] + c \bar{c}$} & $3.2515779884666702 \cdot 10^{-9}\,\mathrm{GeV}^{-4}$ & \multirow{2}{*}{$2.56 \cdot 10^{-14}$} \\
& $3.2515779884665871 \cdot 10^{-9}\,\mathrm{GeV}^{-4}$ & \\ 
\hdashline
 
{$u \bar{u} \to \chi_c\big[{}^3\pwave_{2}^{[1]}\big] + J/\psi\big[{}^3\pwave_{0}^{[8]}\big]$} & 
$1.9155132456571052 \cdot 10^{-14}\,\mathrm{GeV}^{-6}$ & \multirow{2}{*}{$3.29 \cdot 10^{-15}$} \\ 
{\qquad\qquad$ + \eta_c\big[{}^1\swave_{0}^{[1]}\big] + c \bar{c}$}& $1.9155132456570989 \cdot 10^{-14}\,\mathrm{GeV}^{-6}$ & \\ 
\hdashline

{$u \bar{u} \to \chi_c\big[{}^3\pwave_{0}^{[1]}\big] + \chi_c\big[{}^3\pwave_{0}^{[1]}\big]$} & 
$1.8708646257279440 \cdot 10^{-14}\,\mathrm{GeV}^{-4}$ & \multirow{2}{*}{$6.24 \cdot 10^{-15}$} \\ 
{\qquad\qquad$+ \eta_c\big[{}^1\swave_{0}^{[1]}\big] + \eta_c\big[{}^1\swave_{0}^{[1]}\big]$} & 
$1.8708646257279324 \cdot 10^{-14}\,\mathrm{GeV}^{-4}$ & \\ 
\hdashline 

{$g g \to \eta_c\big[{}^1\swave_{0}^{[1]}\big] + J/\psi\big[{}^3\swave_{1}^{[8]}\big] $} & 
$3.3023033137867921 \cdot 10^{-18}\,\mathrm{GeV}^{-4}$ & \multirow{2}{*}{$1.67 \cdot 10^{-14}$} \\ 
{\qquad\qquad$+\Upsilon\big[{}^3\swave_{1}^{[1]}\big] + \eta_c\big[{}^1\pwave_{1}^{[8]}\big]$} & 
$3.3023033137867371 \cdot 10^{-18}\,\mathrm{GeV}^{-4}$ & \\
\hdashline

{$u \bar{u} \to \chi_{c0}\!\big[{}^3\pwave_{0}^{[1]}\big]+ \chi_{b0}\!\big[{}^3\pwave_{0}^{[1]}\big]$} & $1.6923441255992509 \cdot 10^{-25}\,\mathrm{GeV}^{-4}$ & \multirow{2}{*}{$1.74 \cdot 10^{-14}$} \\*
{\qquad\qquad$+\chi_{b0}\!\big[{}^3\pwave_{0}^{[1]}\big]+\eta_c\!\big[{}^1\swave_{0}^{[8]}\big]$} & $1.6923441255992215 \cdot 10^{-25}\,\mathrm{GeV}^{-4}$ & \\

\bottomrule
\end{tabular}
\caption{\it\small Squared-amplitude benchmarks at individual phase-space points. The first column indicates the process. The second column reports the code outputs, with {\tt \mgshort} ({\tt \helaconia}) shown in the top (bottom) row. The last column gives the relative deviation defined in eq.~\eqref{eq:relative_ratio}. 
%The meaning of the {\tt QED} flag is given in the footnote \ref{fnt:QED_flag}.
}
\label{tab:sa_check}
\end{table}\vspace*{-2\baselineskip}
\makeatletter
\renewcommand{\thefootnote}{QED}
\begingroup
  \renewcommand{\@makefnmark}{}%
  \footnotemark
\endgroup
\footnotetext{
This process is generated with the additional {\tt QCD=99 QED=99} flag, which incorporates subdominant contributions by including all diagrams of ${\cal O}(\alpha_s^m\alpha^n)$ for a fixed value of $m + n$. Without this option, \mgshort{} retains only the contributions with the largest non-vanishing power of $\alpha_s$.
\label{fnt:QED_flag}}
\renewcommand*{\thefootnote}{\arabic{footnote}}
\makeatother
\addtocounter{footnote}{-1}
%%%%%%%%%%%%%%%%%%%%%%%%%%%%%%%%%%%%%%%%%%%%%%%%%%%%%%%%%%%%%%%%%%%%%%%%%%%

We have performed this test for almost $1000$ processes. The complete list of benchmarked processes can be found in the supplemental material, while table~\ref{tab:sa_check} presents a selection of representative cases. 
Overall, we observe good agreement between the two codes, with relative deviations typically in the $10^{-14}\textrm{--}10^{-16}$ range, corresponding to double-precision accuracy. Only a small number of processes fail to reach this level of agreement. For these cases, we observe a clear correlation between the magnitude of the matrix element (typically below $10^{-25}$) and the number of coinciding digits between the two codes, with the latter gradually decreasing as the matrix element becomes smaller. This behaviour suggests that the observed discrepancies originate from numerical precision limitations, likely due to cancellations and the interplay between extremely small contributions and finite floating-point precision. 
Importantly, all phenomenologically relevant processes achieve double-precision accuracy.


%%%%%%%%%%%%%%%%%%%%%%%%%%%%%%%%%%%%%%%%%%%%%%%%%%%%%%%%%%%%%%%%%%%%%%%%%%%
\begin{table}[t]
\centering
\renewcommand*{\arraystretch}{1.}
\begin{tabular}[t]{lcc}
\toprule
\multirow{2}{*}{\textbf{process}}& \madgraph & \multirow{2}{*}{\textbf{pull}}  \\
 & \helaconia  \\
\midrule
\multirow{2}{*}{$p p \to g g \to \chi_{b2}\big[{}^3\pwave_{2}^{[1]}\big]$} & $\,{615.100(19)\,\rm nb} $ & \multirow{2}{*}{$0.34$} \\
 & $615.109(19)\,{\rm nb}$ & \\ \hdashline
 \multirow{2}{*}{$p p \to g g \to \eta_{b}\big[{}^1\pwave_{1}^{[8]}\big]$} & $\,{176.4070(54)\,\rm nb} $ & \multirow{2}{*}{$0.26$} \\
 & $176.4048(64)\,{\rm nb}$ & \\ \hdashline
\multirow{2}{*}{$p p \to g g \to \chi_{c0}\big[{}^3\pwave_{0}^{[1]}\big] + H  \ \textsuperscript{\ref{fnt:QED_flag}}$} & $474.64(12)\,{\rm zb} $ & \multirow{2}{*}{$0.40$} \\
 & $474.583(83)\,{\rm zb}$ & \\ \hdashline
\multirow{2}{*}{$p p \to g g \to \chi_{c0}\big[{}^3\pwave_{0}^{[1]}\big] + \eta_b\big[{}^1\swave_{0}^{[1]}\big]$} & $392.63(64)\,{\rm pb}$ & \multirow{2}{*}{$1.21$} \\
 & $393.48(29)\,{\rm pb}$ & \\ \hdashline
 \multirow{2}{*}{$pp \to gg \to \eta_c\big[{}^1\pwave_1^{[8]}\big] + \Upsilon\big[{}^3\pwave_2^{[8]}\big]$} & $18.213(36)\,{\rm pb}$ & \multirow{2}{*}{$0.12$} \\
 & $18.220(49)\,{\rm pb}$ & \\ \hdashline
 \multirow{2}{*}{$pp \to gg \to B^{\ast +}_{c2}\big[{}^3\pwave_2^{[1]}\big] + B^{\ast -}_c\big[{}^3\pwave_1^{[8]}\big]$} & $404.03(16)\,{\rm fb}$ & \multirow{2}{*}{$1.18$} \\
 & $402.92(93)\,{\rm fb}$ & \\ \hdashline
  \multirow{2}{*}{$pp \to u\bar{u} \to h_c\big[{}^1\pwave_1^{[1]}\big] + h_c\big[{}^1\pwave_1^{[1]}\big]$} & $563.32(26)\,{\rm fb}$ & \multirow{2}{*}{$0.81$} \\
 & $564.07(88)\,{\rm fb}$ & \\ \hdashline
 \multirow{2}{*}{$pp \to u\bar{u} \to B_{c1}(H)^+\big[{}^3\pwave_1^{[1]}\big] + b\bar{c}$} & $3.55150(89)\,{\rm pb}$ & \multirow{2}{*}{$1.70$} \\
 & $3.5577(35)\,{\rm pb}$ & \\ \hdashline
  \multirow{2}{*}{$e^+ e^- \to h_b\big[{}^1\pwave_1^{[1]}\big] + Z \ \textsuperscript{\ref{fnt:QED_flag}}$} & $48.2830(80)\,{\rm yb}$ & \multirow{2}{*}{$1.06$} \\
 & $48.2947(76)\,{\rm yb}$ & \\
 \hdashline
  \multirow{2}{*}{$e^+ e^- \to h_b\big[{}^1\pwave_1^{[1]}\big] + ZH \ \textsuperscript{\ref{fnt:QED_flag}}$} & $70.260(21)\,{\rm rb}$ & \multirow{2}{*}{$1.69$} \\
 & $70.314(24)\,{\rm rb}$ & \\ \hdashline
  \multirow{2}{*}{$e^+ e^- \to h_b\big[{}^1\pwave_1^{[1]}\big] + ZHH \ \textsuperscript{\ref{fnt:QED_flag}}$} & $449.23(46)\,{\rm qb}$ & \multirow{2}{*}{$0.53$} \\
 & $448.98(10)\,{\rm qb}$ & \\
\bottomrule
\end{tabular}
\caption{\it\small Cross-section benchmarks for selected $pp$ ($\sqrt{s}=13\,{\rm TeV}$) and $e^+ e^-$ ($\sqrt{s}=1\,{\rm TeV}$) collisions. The first column indicates the process; for $pp$ collisions, only one partonic channel is considered. The second column reports the code outputs, with {\tt \mgshort}~({\tt \helaconia}) shown in the top (bottom). MC uncertainties are given in parentheses. The third column gives the pull defined in eq.~\eqref{eq:pull}.
The meaning of the {\tt QED} flag is given in the footnote \ref{fnt:QED_flag}.}
\label{tab:xsec_check}
\end{table}
%%%%%%%%%%%%%%%%%%%%%%%%%%%%%%%%%%%%%%%%%%%%%%%%%%%%%%%%%%%%%%%%%%%%%%%%%%%

Our next benchmarks test LO integrated cross sections, where we compare the cross sections obtained from eq.~\eqref{eq:mg5_xsec_output} as evaluated in \mgshort\ and \helaconia\ for eleven different processes. The results of this comparison are summarised in table~\ref{tab:xsec_check}.
In this study, we focus on single and double {\pwavetxt} production and employ SM elementary particles to explore more complex phase-space configurations. For proton-proton ($pp$) collisions, we take $\sqrt s = 13\,{\rm TeV}$, while for electron-positron ($e^-e^+$) collisions we use $\sqrt s = 1\,{\rm TeV}$.
To quantify the agreement between the two codes, we define the pull as
\begin{equation}
    \textbf{pull}=\left|\dfrac{\sigma_\mgshort-\sigma_\helaconia}{\sqrt{\Delta_{\mgshort}^2+\Delta_{\helaconia}^2}}\right|,
\label{eq:pull}
\end{equation}
where $\sigma_\mgshort$ and $\sigma_\helaconia$ denote the total cross sections computed with \mgshort\ and \helaconia, respectively, while $\Delta_{\mgshort}$ and $\Delta_{\helaconia}$ denote their corresponding MC uncertainties.
Overall, the results exhibit pulls consistent with statistical fluctuations, demonstrating the stable and reliable performance of the newly implemented \mgshort\ features.

We conclude that the outcomes of our benchmarks provide a robust validation of the extension. This validation includes both the comparison of squared matrix elements at individual phase-space points in \textit{standalone mode} and integrated cross-section comparisons across a representative set of processes. In all cases, we find agreement with the corresponding results obtained with \helaconia, within the expected numerical precision. These tests confirm the reliability of the newly implemented \mgshort\ extension for the automated event generation of processes involving non-relativistic \swave- and \pwavetxt\ bound states.


\subsection{Momentum reshuffling benchmarks}\label{sec:momreshufftest}
To investigate the impact of momentum reshuffling, we study di-bottomonium production with the initial-state and final-state recoil schemes introduced in section~\ref{sec:resh}. In particular, we consider the $2\to2$ process ${gg \to \eta_b[^1 \swave_0^{[1]}]+ \Upsilon(3\swave)[^3 \swave_1^{[8]}]}$ and the $2\to3$ process ${gg \to \eta_b[^1 \swave_0^{[1]}]+\Upsilon(3\swave)[^3 \swave_1^{[8]}]+ Z}$. The large mass difference between the two quarkonium states (${M_{\Upsilon(3\swave)} - M_{\eta_b} = 960\,{\rm MeV}}$) provides an ideal playground for resolving reshuffling effects in an extreme scenario. 
For both processes we adopt two setups. The first configuration is our baseline (see appendix~\ref{app:setup}), with $m_b = 4.7\,\mathrm{GeV}$ corresponding to half of the $\eta_b$ mass ($M_{\eta_b} = 9.4\,\rm GeV$), the mass assigned to the CO state  being $M^{[8]}_{\Upsilon(3\swave)}=M_{\Upsilon(3\swave)}+200\,{\rm MeV}=10.56\,{\rm GeV}$ following the convention used in \pythia{\tt 8}, and the LDMEs taken at their default values as listed in table~\ref{tab:ldme_bottomonia}. In the second configuration, we consider an alternative bottom-mass choice by setting $m_b = M^{[8]}_{\Upsilon(3\swave)}/2 = 5.28\,{\rm GeV}$ and rescaling each LDME according to the mass-dependent factor $\left(5.28/4.7\right)^{3}$ with respect to the first setup, while keeping all other inputs unchanged.

First, we discuss the effect of momentum reshuffling on the squared matrix element as a function of the partonic c.m.\@ energy. The results, shown in the left panel of figure~\ref{fig:etabUps_reshuffle} and in figure~\ref{fig:etabUpsZ_reshuffle}, are obtained by evaluating the squared matrix element at a single phase-space point for each energy value. The phase-space configuration is chosen such that it varies smoothly across the energy range, corresponding to a fixed point in the multidimensional phase-space hypercube. For example, the scattering angles in the $2\to2$ process are kept fixed, while the momenta are rescaled according to the available c.m.\@ energy.

The blue lines are obtained with the default setup ($m_b=4.7\,{\rm GeV}$), employing initial-state reshuffling (dotted), final-state reshuffling with {\tt mom\_resh\_type=1} (dashed), and without reshuffling with $M^{[8]}_{\Upsilon(3\swave)}=2m_b=9.4\,{\rm GeV}$ (solid).
The red lines are instead obtained using $m_b = 5.28\,{\rm GeV}$, combined with initial-state reshuffling (dotted), final-state reshuffling with {\tt mom\_resh\_type=1} (dashed), and without reshuffling with $M_{\eta_b}=2m_b=10.56\,{\rm GeV}$ (solid).



%%%%%%%%%%%%%%%%%%%%%%%%%%%%%%%%%%%%%%%%%%%%%%%%%%%%%%%%%%%%%%%%%%%%%%%%%%%
\begin{figure}[t]
  \centering    
  \input{figures/reshuffling.pgf}
  %\includegraphics{figures/reshuffling.pdf}
  \caption{\it\small Effect of momentum reshuffling on the squared matrix element (left panel) and the cross section (right panel) as functions of the partonic c.m.\@ energy and the final-state invariant mass, respectively, for the process $pp\to gg \to \eta_b[^1 \swave_0^{[1]}]\, \Upsilon(3\swave)[^3 \swave_1^{[8]}]$ at $\sqrt{s}=13\,{\rm TeV}$. Blue (red) lines and bands correspond to the choice $m_b = 4.7\,{\rm GeV}$ ($m_b = 5.28\,{\rm GeV}$). The thick solid lines represent the non-reshuffling case with $M_{\cal Q} = 2 m_b$, while dashed and dotted lines correspond to the final-state and initial-state reshuffling schemes, respectively, with $M_{\eta_b} = 9.4\,{\rm GeV}$ and $M^{[8]}_{\Upsilon(3\swave)} = 10.56\,{\rm GeV}$. The lower panels show ratios with respect to the non-reshuffling case with $m_b=4.7\,\mathrm{GeV}$. The shaded bands indicate the $7$-point scale-variation uncertainties for the scenarios without momentum reshuffling.}
  \label{fig:etabUps_reshuffle}
\end{figure}
%%%%%%%%%%%%%%%%%%%%%%%%%%%%%%%%%%%%%%%%%%%%%%%%%%%%%%%%%%%%%%%%%%%%%%%%%%%



%%%%%%%%%%%%%%%%%%%%%%%%%%%%%%%%%%%%%%%%%%%%%%%%%%%%%%%%%%%%%%%%%%%%%%%%%%%
\begin{figure}[t]
    \centering
    \input{figures/reshuffling2.pgf}
    %\includegraphics{figures/reshuffling2.pdf}
    \caption{\it\small Same as the left panel of figure~\ref{fig:etabUps_reshuffle}, but with an additional $Z$ boson in the final state.}
    \label{fig:etabUpsZ_reshuffle}
\end{figure}
%%%%%%%%%%%%%%%%%%%%%%%%%%%%%%%%%%%%%%%%%%%%%%%%%%%%%%%%%%%%%%%%%%%%%%%%%%%

In the left panel of figure~\ref{fig:etabUps_reshuffle}, we observe that reshuffling effects are relevant only in the low-$\sqrt{\hat s}$ region. The no-reshuffling scenario predicts thresholds corresponding to $\sqrt{\hat s}=4m_b$, which are shifted relative to the physical threshold depending on the chosen value of $m_b$. In contrast, the reshuffled kinematics correctly reproduce the threshold at $\sqrt{\hat s}=M_{\eta_b}+M^{[8]}_{\Upsilon(3\swave)}$, independently of the constituent bottom-quark mass. At higher c.m.\@ energies, the results obtained with and without reshuffling converge, as mass effects become increasingly negligible. This behaviour is expected, since the relative difference between the two approaches scales as $(M_{\cal Q} - 2m_b)/\sqrt{\hat s}$. Thus, the impact of the mass mismatch decreases with increasing energy and vanishes in the asymptotic limit $\sqrt{\hat s}\to\infty$. At the same time, the ratio of the squared matrix elements obtained with the two bottom-quark mass choices approaches the corresponding LDME ratio, namely $(5.28/4.7)^6$.

The final-state reshuffling algorithm reproduces the no-reshuffling results across almost the entire kinematic range, modifying only the threshold behaviour discussed above. This is because in the present $2\to2$ process the reshuffling changes only the magnitudes of the outgoing quarkonium momenta to satisfy the on-shell conditions $M_{\cal Q}=2m_b$ for both $\eta_b$ and $\Upsilon(3\swave)$, while leaving the scattering angles and the partonic c.m.\@ energy unchanged. As a result, all Mandelstam invariants are preserved, and the Lorentz-invariant matrix element remains identical to that obtained without reshuffling.
By contrast, the initial-state recoil algorithm enforces the on-shell condition by modifying the momenta of the incoming partons. 
For $m_b=4.7\,\mathrm{GeV}$, the partonic c.m. energy compensates for the replacement of the (larger) $\Upsilon(3\swave)$ physical mass by $2m_b$.  
As a result, the matrix element corresponding to a physical collision energy $\sqrt{\hat s}$ is evaluated at a lower effective partonic c.m.\@ energy, shifting the spectrum towards larger $\sqrt{\hat{s}}$ values. For $m_b=5.28\,\mathrm{GeV}$, the opposite effect occurs, as the replacement of the (lower) $\eta_b$ physical mass by $2m_b$ requires an increase in the collision energy to compensate for the excess. The matrix element is therefore evaluated at a higher effective c.m.\@ energy, shifting the spectrum towards smaller $\sqrt{\hat{s}}$ values.

The $2\to3$ process shown in figure~\ref{fig:etabUpsZ_reshuffle} exhibits a more intricate dependence of the squared matrix element on the treatment of the bound-state masses. As in the previous case, the reshuffling methods reproduce the physical threshold, whereas the approaches without reshuffling exhibit a systematic shift depending on the bottom-quark mass. While this behaviour is expected, the structure of the matrix element in the threshold region becomes much richer and less predictable. This originates from the non-trivial dependence of the Mandelstam variables on the chosen reshuffling algorithm, leading to the complex shapes observed. These effects are particularly pronounced for the final-state reshuffling method, which deviates significantly from the non-reshuffling case.
In the high-energy limit $\sqrt{\hat{s}}\to\infty$, all methods again converge to the same asymptotic behaviour, as mass effects become increasingly negligible.

To extend the analysis beyond the matrix-element level and assess the phenomenological impact of reshuffling effects, we also study the invariant-mass spectrum of the process $gg \to \eta_b[^1 \swave_0^{[1]}] + \Upsilon(3\swave)[^3 \swave_1^{[8]}]$, shown in the right panel of figure~\ref{fig:etabUps_reshuffle}. At LO, the invariant mass $M_{\mathcal{QQ}}$ of the final-state $\eta_b$--$\Upsilon(3\swave)$ system coincides with the partonic c.m.\@ energy, $M_{\mathcal{QQ}}=\sqrt{\hat s}$. The spectrum therefore provides a direct connection between the squared matrix elements shown in the left panel and their corresponding contribution to the differential cross section.
The cross section, however, is not only affected by modifications of the matrix element induced by the reshuffling procedure, but also by changes in the available phase-space volume. For $m_b=4.7\,\mathrm{GeV}$, the physical meson masses exceed the constituent-mass values, leaving less kinetic energy available to populate the phase space. This effect suppresses the cross section in the threshold region and shifts the reshuffled spectra towards larger $M_{\mathcal{QQ}}$ values compared to the non-reshuffling case. For $m_b=5.28\,\mathrm{GeV}$, the opposite behaviour is observed. In this case, the physical threshold lies below the corresponding threshold at $4m_b$, such that the excess energy increases the accessible phase-space volume and enhances the cross section near threshold. The resulting spectra are therefore shifted towards smaller values of $M_{\mathcal{QQ}}$.
Since the final-state reshuffling algorithm leaves the matrix element essentially unchanged, the observed deviations in the corresponding cross sections originate almost entirely from the modified phase-space treatment. By contrast, the initial-state recoil prescription affects both the matrix element and the phase-space measure. Since these effects act in the same direction, they lead to a significantly larger deviation from the non-reshuffling prediction than that observed at the matrix-element level alone.

The more realistic phase-space treatment, in which the measured bound-state masses determine the available phase-space volume, suggests that applying a momentum-reshuffling prescription can significantly improve the physical description of the predictions obtained with \mgshort\ compared to calculations performed without reshuffling. However, we observe sizeable differences between the various reshuffling prescriptions. These differences reflect an intrinsic ambiguity in the mapping between the physical bound-state kinematics and the constituent-level kinematics required for the matrix-element evaluation. Thus, they should be regarded as an additional source of theoretical uncertainty.

